\chapter{The review: what the benchmark can and cannot say}\label{ch:review}

On 10 November the whole project was reviewed critically: every document, script, test and saved run. The result is
\code{FABLE_REVIEW.md}, about 800 lines. This chapter teaches what is in it, because it changed how every number in
this booklet should be read. The next four chapters are the experiments that the review asked for.

\section{How the review was done}

Nothing was trained. The review read the files, and it ran small NumPy checks on the saved predictions
(\code{scripts/review_checks.py}, about two minutes). Every statement carries a tag that says how far to trust it:
\begin{itemize}
\item \textbf{[file]}: a number read from a repository file or a saved run.
\item \textbf{[check]}: a number the review computed itself from saved arrays. These used the two COMSOL configurations
that are on the Mac, so they are diagnostics on test data, not results.
\item \textbf{[web]} and \textbf{[memory]}: a reference found by a search, or known but not fetched again.
\item \textbf{[not verified]}: stated in our documents, and not confirmed.
\end{itemize}

\section{The verdict}

\begin{keyidea}[The review in three sentences]
The project asks the right question, and its last move, a network that works in the current domain, was the right one.
But most of the effort had gone into ranking field-domain networks by 1 to 3 points on a benchmark that cannot support
such rankings. The benchmark has two wire shapes, one seed per method, a noisy target and an unexplained physics
anomaly, and it had been used to compare more than twenty method variants.
\end{keyidea}

\section{What went right}

It is worth knowing these too, because they are what you can defend without hesitation.
\begin{itemize}
\item \textbf{Test hygiene.} Since the restart, COMSOL is kept out of training, normalisation and model selection, and
the code enforces it.
\item \textbf{Faithful baselines.} The MAGIC-UNet port is checked layer by layer against the released model, and the
run 1 failure was diagnosed, not tuned away.
\item \textbf{Asking for the limit.} The Wiener oracle and the lost-band floor (Chapter~\ref{ch:floor}) are the right kind
of question.
\item \textbf{The ambiguity analysis} of Chapter~\ref{ch:question} is correct, and it is the core of a possible paper.
\item \textbf{The run 3 diagnosis} is a real finding, not a bug report.
\end{itemize}

\section{What the ``15 pairs'' really are}

Every score in this booklet is a mean over ``15 COMSOL pairs''. Look at what they are:

\begin{center}
\small
\begin{tabularx}{\textwidth}{>{\raggedright\arraybackslash}p{4.4cm}X}
\toprule
What it looks like & What it is \\
\midrule
15 test images & 5 coil configurations $\times$ 3 field components \\
5 configurations & 2 wire shapes. \code{C1_C2}, \code{C1_C2_C3} and \code{C1_C2_C3_C4} are mirror copies of
\code{L1_C1} (correlation 0.97--0.98) \\
Independent noise in each & The same mesh noise: all four C1 configurations come from one COMSOL solve \\
Unseen data & The simulator was built around these coils: its layout families, wire widths, current and noise level
all come from them \\
One look per method & At least 22 method variants scored on the same 15 pairs. FISTA's $\lambda$ was retuned after
network scores were known; FISTA-Net got a second run because of its first COMSOL score \\
\bottomrule
\end{tabularx}
\end{center}

\begin{careful}[What follows from this table]
Each step was defensible on its own. Together they mean two things. First, the test set measures performance on the two
shapes the simulator was designed for, not on new geometries. Second, when many variants are tried on one small test
set, the best one looks better than it is. A difference of 1 to 3 points, on one seed, is not a finding.
\end{careful}

\section{What went wrong, sorted by cause}

The review sorts the problems into five kinds. This is the useful part: it tells you \emph{which kind} of fix each one
needs.

\textbf{(a) Information limits.} Part of the 0.5\um\ field is simply not in the 3\um\ measurement. This is real, but the
``floor'' of 34.5\% in Chapter~\ref{ch:floor} overstates it, because it mixes two different things: real field that was
lost, and mesh noise in the target that no method could ever reproduce. Chapter~\ref{ch:limits} separates them.

\textbf{(b) Modelling errors.} The blur model $e^{-k\,\Delta z}$ is physics, but COMSOL does not follow it exactly. At
3\um\ the mismatch is 10--17\%. At 10\um\ it is 64--140\%, more than any real field could decay. Chapter~\ref{ch:comsol}
finds out why. COMSOL's three components are also not consistent with each other, and MAGIC-UNet's two outputs $J_x$ and
$J_y$ come from separate decoders, so its current does not conserve charge exactly.

\textbf{(c) Data problems.} Two test shapes, shared mesh noise, a simulator built around them, no sensor noise anywhere,
and no real measurement.

\textbf{(d) Optimisation and architecture.}
\begin{itemize}
\item \emph{CSC-LISTA is under-trained, not data-limited.} At 64 layers it has 49 million parameters, trains for 1{,}000
steps and keeps the checkpoint from step 100--300. On its \emph{own} training distribution it loses to plain FISTA at the
same depth with the same dictionary (48.8\% against 45.3\% at 64). A learned network that cannot beat the method it
starts from has an optimisation problem.
\item \emph{It is not doing sparse coding.} 99\% of its code values are non-zero. With dense codes it behaves almost like a
linear filter, which is why it ends up next to the best linear filter (Chapter~\ref{ch:floor}).
\item \emph{FISTA-300 is not converged.} Its error was still falling from 32 to 64 to 300 iterations, and nobody had run
1{,}000. Chapter~\ref{ch:fista_depth}.
\item \emph{MAGIC-UNet has no data-consistency step and a small view.} Nothing forces its output to agree with the
measurement, and its 12.8\um\ tile is only about four standoffs wide. Chapters~\ref{ch:consistency} and~\ref{ch:ten}.
\end{itemize}

\textbf{(e) Evaluation.} The metric is taken against a noisy target, with one seed and no error bars, and models are
selected on 64 crops.

\section{The surprising results, one by one}

For each result that had seemed strange, the review gives the explanation that the evidence supports best, and a cheap
test that could prove that explanation wrong. Most of these tests were then run.

\begin{center}
\small
\begin{tabularx}{\textwidth}{>{\raggedright\arraybackslash}p{3.3cm}>{\raggedright\arraybackslash}X>{\raggedright\arraybackslash}p{4.6cm}}
\toprule
Result & Best-supported explanation & What the test then showed \\
\midrule
CSC-LISTA got worse with 8{,}000 layouts (33.8 to 36.0\%) & Checkpoint and seed noise. Both runs kept a checkpoint from
step 100--200, before the extra layouts could matter & Not run. The claim ``data volume is not the limit'' of
Chapter~\ref{ch:followups} is not supported \\
Learned field-domain networks sit next to the Wiener oracle & With dense codes they act almost linearly, so they can only
restore the measured band & Holds: they leave 79--82\% of the next band unrecovered (Chapter~\ref{ch:limits}) \\
FISTA-300 beats every learned field-domain network & FISTA was still improving at 300 iterations; the networks stop at 64
layers & FISTA is best at 1{,}000 iterations: 26.7\%, not 30.5\% (Chapter~\ref{ch:fista_depth}) \\
MAGIC-UNet beats methods that were given the standoff & Its current-domain prior fills in frequencies just above what
the measurement carries & Against FISTA-1000 the raw network only ties; with one correction it wins
(Chapter~\ref{ch:consistency}) \\
Run 3 fails on the IC coil & Each tile decides its own height & One global height does not rescue it: 52.0 to 46.9\%
(Chapter~\ref{ch:height}) \\
\bottomrule
\end{tabularx}
\end{center}

\begin{keyidea}[Notice the last row]
The review's own prediction was wrong there, and the document says so. A good review states its explanations in a
form that an experiment can refute, and then reports the refutation.
\end{keyidea}

\section{Corrections to earlier chapters of this booklet}

\begin{careful}[Read these before quoting Chapters 9 to 19]
\begin{itemize}
\item \textbf{Chapter~\ref{ch:floor}:} ``any method below 34.5\% must be filling in lost frequencies'' is too strong. The
floor is inflated by noise in the target. See Chapter~\ref{ch:limits}.
\item \textbf{Chapter~\ref{ch:followups}:} ``data volume is not what limits CSC-LISTA'' and ``the limit is the simulator
gap'' are not supported by that experiment.
\item \textbf{Chapter~\ref{ch:defence}:} ``$h$ stays fixed because it is physics and we know it exactly'' is true of physics.
It is not true of this COMSOL export, which does not obey free-space physics (Chapter~\ref{ch:comsol}).
\item \textbf{Chapter~\ref{ch:question}:} ``the height cannot come from the field alone'' is right, but the ambiguity is
\emph{one-sided}. Without any prior the data still give an upper bound on the depth (Chapter~\ref{ch:limits}).
\item \textbf{Chapters~\ref{ch:results} to~\ref{ch:magic_results}:} every comparison with ``FISTA'' there means FISTA at 300
iterations, which was stopped too early.
\end{itemize}
\end{careful}

\section{What the review said to stop, and to do}

\textbf{Stop:} new CSC-LISTA, LFISTA and FISTA-Net variants; reading differences under about 3 points; reporting
5--10\um\ COMSOL scores as physics before the decay anomaly is explained; treating $B_x$, $B_y$ and $B_z$ as three
unrelated images; quoting the lost-band floor as a limit of the problem.

\textbf{Do}, as a list of milestones:

\begin{center}
\small
\begin{tabularx}{\textwidth}{l>{\raggedright\arraybackslash}X>{\raggedright\arraybackslash}p{4.2cm}}
\toprule
 & Milestone & Status on 16 November \\
\midrule
M1 & Re-score run 2 and run 3 with a data-consistency step and a height read-out. No training & Done:
Chapters~\ref{ch:height} and~\ref{ch:consistency} \\
M2 & Repair the measuring stick: noise edge, scores by band, FISTA run longer, more seeds & Done:
Chapters~\ref{ch:limits}, \ref{ch:fista_depth}, \ref{ch:consistency} \\
M3 & Open the COMSOL model and explain the fast decay; export new geometries & First half done:
Chapter~\ref{ch:comsol}. No new geometries \\
M4 & Classical test of whether the height can be identified & Not done \\
M5--M6 & A new solver: learned current prior inside data-consistency steps, with a search over height & Not done \\
M7--M9 & Generalisation tests, real QDM data, uncertainty & Not done \\
\bottomrule
\end{tabularx}
\end{center}

What the review recommends building after these milestones is in Chapter~\ref{ch:next}.

\section{What is already published, and what is left to claim}

The review searched the literature so that you do not claim something that exists.

\textbf{Do not claim as new:} learning the map from field to current (Reed et al.); using all three field components
(Broadway et al.\ 2020); a sparse or total-variation prior on the current; fitting the sensor distance jointly with
the source as an idea (Setescak et al.\ 2026, for magnetisation); unrolled networks that are told the blur (USRNet);
``a network that works without being told the standoff'', \emph{if} Reed et al.\ already trained over $50\pm10$\um.
The review marks that point as not verified, and Chapter~\ref{ch:magic} explains why it is open.

\textbf{Still open, if the experiments deliver:} inferring the standoff for the whole image with a learned current prior
and exact physics; the one-sided analysis of the ambiguity with a measured height resolution; a reconstruction that
states which frequencies came from the data and which from the prior; validation across standoffs and on a real QDM.

\begin{defend}[If your professor asks: was most of the earlier work wasted?]
No, but it was aimed at the wrong target. The unrolled networks taught us where the limit of field-domain methods is,
and the review used their saved outputs to show why. What was wasted was the ranking: differences of a few points
between them were never real.
\end{defend}

\begin{defend}[If your professor asks: how do you know your winner is not just the luckiest of 22 tries?]
You do not, from the old test alone. That is why the next step was a protocol written before any new score
(Chapter~\ref{ch:fista_depth}), three seeds, and a rule fixed in advance for what counts as a win
(Chapter~\ref{ch:consistency}).
\end{defend}
